% 原文 pp.203--206
\subsubsectionlink{量子Brown運動模型とCaldeira--Leggett模型の限界}
5.2.4節では、Caldeira--Leggett模型から、空間的デコヒーレンス率の見積もりについて、非常に簡単な式(5.61)が得られることを示した。特に、空間的デコヒーレンス率と散逸率の比は、コヒーレントな隔たりと熱的de Broglie波長の比によって与えられた。2.11節で示したように、中間的あるいは巨視的な大きさの物体が、同じく中間的・巨視的な距離だけ離れた二つの位置の重ね合わせで記述される、典型的状況では、この比は非常に大きくなりうる。これは、そのような状況では、典型的なデコヒーレンス率が、散逸率より何桁も大きくなることを示唆する。

しかし、このような見積もりを、あまり文字どおりに解釈しないよう注意すべきである。Caldeira--Leggett方程式(5.57)、したがってデコヒーレンス率の表式(5.61)を導いた仮定は、物理的に現実的で興味ある場合には、しばしば満たされない\cite{ref162,ref228,ref229}。したがって、(5.61)を、空間的重ね合わせのデコヒーレンス率について普遍的に成り立つ表式と解釈し、過度に一般化するのは適切でない。さらに、Caldeira--Leggett方程式(5.57)から生じる、デコヒーレンスと散逸の一般的な動力学は、この模型の文脈と範囲の中で考えるべきであり、この模型から得た結果から、過度に一般的な結論を引き出すことには慎重であるべきだ。

例えば、(5.61)のデコヒーレンス率は、位置空間の重ね合わせ状態のコヒーレントな隔たり$\Delta X=|X-X'|$の二乗に比例する。もちろん、隔たりが大きいほど重ね合わせは非古典的であり、デコヒーレンスも速くなると期待されるので、直観的には妥当である。

% 原文 p.204
\hypertarget{srcpage:204}{}\label{srcpage:204}
しかし、物理的には、このデコヒーレンス率には、隔たりをさらに増しても超えられない上限があるはずである。

この上限の理由は、第3章で論じた環境散乱の結果を思い出せば容易にわかる。そこでは、散乱される環境粒子の典型的波長が、コヒーレントな隔たり$\Delta X$に比べて短い「短波長極限」（3.3.1節参照）と、逆に長い「長波長極限」（3.3.2節参照）を区別した。短波長極限では、一回の散乱で隔たり$\Delta X$を完全に識別できるが、長波長極限では、顕著な量の経路情報を環境に伝えるまでに、多数の事象が必要である。3.3.2節で論じたように、短波長極限で得られるデコヒーレンス率は、そのため上限を与える。環境はそれ以上の経路情報を得られないので、デコヒーレンス率は、短波長極限で計算した値を超えられない。

同様の議論は、ここでの量子Brown運動模型にも成り立つと期待される。すなわち、調和振動子の環境は、ある最大のコヒーレンス長をもつはずである。隔たり$\Delta X$がその長さを超えると、環境散乱で考えた短波長領域に相当する領域に入り、デコヒーレンス率が飽和すると予想される。しかし明らかに、Caldeira--Leggett模型では、このような飽和は起こらない。環境のコヒーレンス長に相当する量が模型に入る場所はなく、隔たり$\Delta X$を増すと、デコヒーレンス率(5.61)は際限なく増加する。

この例は、Caldeira--Leggett模型の結果を、任意の物理的状況とパラメータへ外挿できないことを、はっきり示している。この点をさらに調べるため、質量をもつ粒子が、質量のないスカラー場に結合する一般的な模型を考えられる。この模型はUnruhとZurek\cite{ref109}が詳しく研究した。環境場の支配的な波長が、中心粒子の典型的なコヒーレントな隔たり$\Delta X$より、はるかに長いと仮定すると（双極子型の近似に相当する）、この粒子・場模型の極限として、量子Brown運動模型が出現することが示された\cite{ref015,ref109}。

これは、Caldeira--Leggett模型だけでなく、マスター方程式(5.25)で記述される、より一般的な模型も含めて、量子Brown運動模型の基礎に、暗黙の長波長仮定があることを示す。これによって、Caldeira--Leggett模型の予測するデコヒーレンス率が飽和しない理由も説明される。実際、一般的な粒子・場模型では、$\Delta X$が場の支配的波長より小さい場合、デコヒーレンス率は$\Delta X$にほぼ二次的に依存し、$\Delta X$がはるかに大きくなると一定値に近づくことがわかっている\cite{ref015,ref109,ref229}。

% 原文 p.205
\hypertarget{srcpage:205}{}\label{srcpage:205}
したがって、(5.61)に明らかな、デコヒーレンス率の飽和の欠如は、基礎にある模型の人為的な帰結にすぎない\cite{ref228,ref229}。

量子Brown運動のMarkov模型のもう一つの限界は、時間に依存しない係数(5.26)の使用にある。5.2.2節の最後で述べたように、時間依存係数を用いた、より現実的な扱い［(5.42)参照］では、通常拡散係数$D$は、最終値(5.48)へ落ち着く前に、初期値ゼロから増大する、温度に依存しない過渡応答を示す。その時間尺度は高振動数カットオフ$\Lambda$［(5.44)参照］で定まり、したがって$\Lambda$はデコヒーレンス率に限界を課す。この扱いでは、その過渡応答を完全に無視し、代わりに係数$D$は常に一定と仮定した。

さらに、初め系と環境の状態が完全に無相関であるという仮定(4.23)が、かなり見かけ上のデコヒーレンスの動力学を誘起する役割については、AnglinとZurek\cite{ref229}が詳しく論じている。最後に、Caldeira--Leggett模型は高温環境の仮定に基づくため、例えば超伝導体を含むような低温実験でのデコヒーレンスの表現には、不十分かもしれない。特に低温では、完全なマスター方程式(5.25)の異常拡散項(5.40)が重要になり、通常拡散項(5.34)と同程度の大きさに近づいて、系の運動に大きく影響することがある\cite{ref015,ref220}。

したがって、特定のデコヒーレンス模型から得た結論を、自然界のデコヒーレンスの一般的な性質について、普遍的に有効な言明だと誤解してはならない。見てきたように、デコヒーレンス模型に現れる、一見物理的に重要そうな性質も、しばしば、模型とその導出の仮定の人為的な帰結にすぎないことがある。

ただし、こうした注意にもかかわらず、Caldeira--Leggett模型から得るデコヒーレンス率の見積もりが、驚くほど頑健であることも、述べておかなければならない。すなわち、その見積もりを、より一般的で現実的な模型から導いた予測と比較すると、より大きな食い違いが予想される場合でさえ、よい一致が見られる。例えば、Caldeira--Leggett模型の基礎にある高温仮定が成立しない環境でも、この、厳密にいえばもはや誤った模型から得られるデコヒーレンス率は、より複雑な非Markov模型の値に、しばしば驚くほど近い（さらに詳しい比較と議論は、Paz、Habib、Zurekの解析\cite{ref184}参照）。この性質と、見積もりを得る簡単さ［(5.61)参照］により、Caldeira--Leggett模型は、多くのデコヒーレンス問題を扱う、優れた出発点となる。

% 原文 p.206
\hypertarget{srcpage:206}{}\label{srcpage:206}
\subsubsectionlink{厳密なマスター方程式}
線形量子Brown運動模型には、厳密に解けるという顕著な性質がある。おそらくさらに驚くべきことに、他の著者の先行研究\cite{ref109,ref221,ref230,ref231,ref232}を受けて、Hu、Paz、Zhangが1992年に初めて示したように\cite{ref219}、得られる厳密な非Markovマスター方程式は、任意の環境スペクトル密度に対して時間局所的である（4.4節の議論も参照）。このマスター方程式は、上の5.2.1節、5.2.2節で導いたBorn--Markovマスター方程式(5.25)と、まったく同じ関数形を取る。すなわち
\chfiveeq{84}
であり、$\hat H_{\mathrm S}=\hat P^2/(2M)+M\Omega^2\hat X^2/2$は、系の調和振動子の自己Hamiltonian(5.23)を表す。

厳密な方程式(5.84)とBorn--Markov版(5.25)の唯一の、そして決定的な違いは、係数$\widetilde\Omega^2$、$\gamma$、$D$、$f$が、ここでは陽に時間に依存し、他の時間依存係数の、かなり複雑な関数となることである。それらの係数は、雑音核と散逸核(5.16)、(5.17)のさまざまな積分を含む（詳細は文献\cite{ref219}参照）。厳密な非Markov動力学の微妙な特徴は、すべて(5.84)の係数$\widetilde\Omega^2(t)$、$\gamma(t)$、$D(t)$、$f(t)$にまとめられる一方、右辺の各項の解釈は、Born--Markovマスター方程式(5.25)とまったく同じである。4.3節で簡単に述べたように、厳密なマスター方程式を、正値性を保証するLindblad形(4.44)で書けないこと自体は意外でないが、それにもかかわらず、これらの係数の特定の時間依存性が、縮約密度行列の正値性［(4.42)参照］を保証する。

線形量子Brown運動の模型で、縮約された運動を厳密に決定でき、しかも時間局所的なマスター方程式で与えられることは、かなり特異な性質である。物理的に興味ある他の模型で、厳密なマスター方程式を計算できるものは、あまり多くないように見える。厳密解を許す一例は、特定の形の相互作用Hamiltonianで記述する、調和振動子環境と二準位系の相互作用である（これはスピン・ボソン模型の一例である。次の5.3節参照）。Garraway\cite{ref233,ref234}が示したように、得られる厳密なマスター方程式は、量子Brown運動のマスター方程式(5.84)と同様、時間局所的で、時間依存係数をもつ。
